\documentclass{beamer}

\usepackage[T2A]{fontenc}
\usepackage[utf8]{inputenc}
\usepackage[english,russian]{babel}
\input{settings.tex}


\title{Псевдообращение}
\subtitle{Вычислительный интеллект, лекция 2}
\author{Сергей Савин}
\centering
\date{\mydate}



\begin{document}
\maketitle





\begin{frame}{Метод наименьших квадратов}
	\begin{flushleft}
		
		Рассмотрим следующую задачу: найти $\bo{x}$, минимизирующий $|| \bo{A}\bo{x} - \bo{y} ||_2^2$. Это \emph{задача наименьших квадратов}. 
		
		\bigskip
		
		\begin{itemize}
			\item Предполагаем, что $\bo{A} \in \R^{m \times n}$ имеет линейно независимые столбцы, $m \geq n$.
			
			\item Величина $\bo{e} = \bo{A}\bo{x} - \bo{y}$ называется невязкой.
			
			\item Задача наименьших квадратов заключается в поиске решения с \emph{наименьшей невязкой}.
		\end{itemize}
		
		Заметим, что $|| \bo{A}\bo{x} - \bo{y} ||_2^2 = (\bo{A}\bo{x} - \bo{y})\T(\bo{A}\bo{x} - \bo{y}) = \bo{x}\T\bo{A}\T\bo{A}\bo{x}
		-\bo{x}\T\bo{A}\T\bo{y} -\bo{y}\T\bo{A}\bo{x} +  \bo{y}\T \bo{y}$.
		
		
	\end{flushleft}
\end{frame}


\begin{frame}{Метод наименьших квадратов $\rightarrow$ левая обратная}
	\begin{flushleft}
		
		Найдём экстремум:
		%
		\begin{equation}
			\frac{\partial}{\partial \bo{x}}  \left( 
			\bo{x}\T\bo{A}\T\bo{A}\bo{x}
			-\bo{x}\T\bo{A}\T\bo{y} -\bo{y}\T\bo{A}\bo{x} +  \bo{y}\T \bo{y}
			\right) 
			= 0
		\end{equation}
		\begin{equation}
			\bo{x}\T\bo{A}\T\bo{A} + \bo{x}\T\bo{A}\T\bo{A} - 2\bo{y}\T\bo{A} 
			= 0
		\end{equation}
		\begin{equation}
			2\bo{A}\T\bo{A}\bo{x} - 2\bo{A}\T\bo{y} 
			= 0
		\end{equation}
		\begin{equation}
			\bo{A}\T \bo{A}\bo{x} = \bo{A}\T\bo{y}
		\end{equation}
		\begin{equation}
			\bo{x} = (\bo{A}\T \bo{A})^{-1} \bo{A}^\top\bo{y}
		\end{equation}
		
		Определим \emph{левую обратную матрицу}:
		%
		\begin{equation}
			\bo{A}^+_l = (\bo{A}\T \bo{A})^{-1} \bo{A}\T
		\end{equation}
		
		
	\end{flushleft}
\end{frame}



\begin{frame}{Левая и правая обратные матрицы}
	\begin{flushleft}
		
		Заметим, что наш вывод предполагает, что $\bo{A}\T \bo{A}$ — матрица полного ранга, а значит $\bo{A}$ имеет полный столбцовый ранг.
		
		\bigskip
		
		Псевдообратная матрица $\bo{A}^+_l = (\bo{A}\T \bo{A})^{-1} \bo{A}\T$ называется "левой обратной" в силу следующего тождества:
		%
		\begin{align}
			\bo{A}^+_l  \bo{A} =  \bo{I}
			\\
			(  \textcolor{myblue} {\bo{A}\T\bo{A}})^{-1} \textcolor{myblue} {\bo{A}\T\bo{A}} =  \bo{I}
		\end{align}
		
		\bigskip
		
		Аналогично, правая обратная определяется как $\bo{A}^+_r = \bo{A}\T (\bo{A}\bo{A}\T )^{-1}$:
		%
		\begin{align}
			\bo{A} \bo{A}^+_r=  \bo{I}
			\\
			\textcolor{myblue}{\bo{A}\bo{A}\T} (\textcolor{myblue}{\bo{A}\bo{A}\T} )^{-1} =  \bo{I}
		\end{align}
		
	\end{flushleft}
\end{frame}


\begin{frame}{Псевдообратная матрица}
	\begin{flushleft}
		
		{\small
			\begin{itemize}
				
				\item Если матрица $\bo{A}$ имеет \emph{полный столбцовый} ранг, то её левая обратная равна её \emph{псевдообратной}: $\bo{A}^+_l = \bo{A}^+$.
				
				\item Если матрица $\bo{A}$ имеет \emph{полный строчный} ранг, то её правая обратная равна её псевдообратной: $\bo{A}^+_r = \bo{A}^+$.
				
			\end{itemize}
			
			Для матрицы $\bo{A}$ с SVD разложением $\bo{A} = \bo{U} \Sigma \bo{V}'$ псевдообратная матрица определяется как:
			
			\begin{equation}
				\bo{A}^+ = \bo{V} \Sigma^i \bo{U}'
			\end{equation}
			
			где
			
			\begin{align*}
				\Sigma = \begin{bmatrix}
					\sigma_1 & 0 & ... & 0 & 0 & .... \\
					0 & \sigma_2 & ... & 0 & 0 & .... \\
					... & ... &  ... & ... & ... & .... \\
					0& 0& ...& 0 & 0& ... \\
					0& 0& ...& 0 & 0& ... \\
					... & ... &  ... & ... & ... & .... 
				\end{bmatrix} 
				\ \ \
				\Sigma^i  = \begin{bmatrix}
					1 /  \sigma_1 & 0 & ... & 0 & 0 & .... \\
					0 & 1 /  \sigma_2 & ... & 0 & 0 & .... \\
					... & ... &  ... & ... & ... & .... \\
					0& 0& ...& 0 & 0& ... \\
					0& 0& ...& 0 & 0& ... \\
					... & ... &  ... & ... & ... & .... 
				\end{bmatrix}
			\end{align*}		
		}
		
	\end{flushleft}
\end{frame}







\begin{frame}{Вычисление псевдообратной матрицы}
	\begin{flushleft}
		
		Псевдообратную матрицу легко вычислить. Мы можем вызвать функцию \texttt{pinv} в MATLAB или Python:
		%
		
		\begin{align*}
			\bo{A}^+ \texttt{= pinv}(\bo{A})
		\end{align*}
		
		\bigskip
		
		Функция \texttt{pinv} вычисляет псевдообратную матрицу на основе SVD разложения. Вычисление левой (или правой) обратной матрицы на основе определения является вычислительно затратным и нестабильным.
		
		\bigskip
		
		Часто решение линейной системы $\bo{A}\bo{x} = \bo{y}$ напрямую происходит гораздо быстрее, чем вычисление $\bo{x} = \bo{A}^+\bo{y}$:
		
		%double-check
		\begin{align*}
			\texttt{x = A \textbackslash \  y}
		\end{align*}
		
		
	\end{flushleft}
\end{frame}






\begin{frame}{Ортонормированные матрицы}
	\begin{flushleft}
		
		Пусть матрица $\bo{M}$ является ортонормированной (но не обязательно квадратной), то есть $\bo{M}\T \bo{M} = \bo{I}$. Её псевдообратная матрица может быть упрощена:
		
		\begin{align}
			\bo{M}^+ = (\bo{M}^\top\bo{M})^{-1} \bo{M}^\top = \bo{M}\T
		\end{align}
		
		\bigskip
		
		Тогда решение методом наименьших квадратов для уравнения $\bo{M} \bo{x} = \bo{y}$ может быть найдено как:
		%
		\begin{align}
			\bo{x}_{LS} = \bo{M}^\top \bo{y}
		\end{align}
		
		Если $\bo{M}$ ортонормирована и квадратна, то $\bo{M}\T = \bo{M}^{-1} = \bo{M}^+$.
		
		
	\end{flushleft}
\end{frame}




\begin{frame}{Вычисление невязки}
	\begin{flushleft}
		
		Для уравнения $\bo{A}\bo{x} = \bo{y}$ и решения методом наименьших квадратов $\bo{x}_{LS} = \bo{A}^+\bo{y}$ вычислим невязку $\bo{e} = \bo{y} - \bo{A}\bo{x}_{LS}$. Подставим решение:
		
		\begin{equation}
			\bo{e} = \bo{y} - \bo{A}\bo{A}^+\bo{y}
		\end{equation}
		
		Заметим, что:
		
		\begin{itemize}
			\item Невязка может быть найдена как $\bo{e} = (\bo{I} - \bo{A}\bo{A}^+)\bo{y}$.
			
			\item Ближайшее значение $\bo{A}\bo{x}$ к $\bo{y}$ равно $\bo{y}^* = \bo{A}\bo{A}^+\bo{y}$.
			
		\end{itemize}
		
	\end{flushleft}
\end{frame}




\begin{frame}{Нулевое пространство, 1}
	\begin{flushleft}
		{\small
		Псевдообратная матрица — это инструмент для нахождения решения системы $\bo{A}\bo{x} = \bo{y}$ с \emph{наименьшей нормой и наименьшей невязкой}. Чтобы найти все остальные решения, нам необходимо изучить \emph{нулевое пространство} (ядро) матрицы $\bo{A}$.
		
		Найдём все решения следующей задачи:
		%
		\begin{equation}
			\begin{bmatrix}
				1 & 2 & 0 \\ 0 & 3 & 0
			\end{bmatrix}
			\begin{bmatrix}
				x_1 \\ x_2 \\ x_3
			\end{bmatrix}
			=
			\begin{bmatrix}
				0 \\ 0
			\end{bmatrix}
		\end{equation}
		%
		Любой $\bo{x}$, для которого $x_1 = 0$ и $x_2 = 0$, удовлетворяет этому уравнению. Таким образом, все решения этой задачи имеют вид:
		%
		\begin{equation}
			\begin{bmatrix}
				x_1 \\ x_2 \\ x_3
			\end{bmatrix}
			=
			\begin{bmatrix}
				0 \\ 0 \\ z
			\end{bmatrix}
			=
			\begin{bmatrix}
				0 \\ 0 \\ 1
			\end{bmatrix}
			z, \ \forall z
		\end{equation}
		
		где $\begin{bmatrix}
			0 & 0 & 1
		\end{bmatrix}\T$ является \emph{базисом нулевого пространства} для матрицы $\bo{A}$ этой системы, а $z$ — координата в нулевом пространстве.
		}
	\end{flushleft}
\end{frame}



\begin{frame}{Нулевое пространство, 2}
	\begin{flushleft}
		
		\emph{Нулевое пространство} (ядро) матрицы $\bo{A}$ — это множество всех векторов, которые матрица отображает в ноль. Эквивалентно, это множество всех решений уравнения $\bo{A}\bo{x} = 0$.
		
		\bigskip
		
		Мы можем найти ортонормированный базис в нулевом пространстве $\bo{A}$, вызвав функцию \texttt{null}:
		
		\begin{align*}
			\texttt{N = null(A)}
		\end{align*} 
		
		По определению, $\bo{A}\bo{N} = 0$.
		
	\end{flushleft}
\end{frame}





\begin{frame}{Все решения}
	\begin{flushleft}
		
		Если система уравнений $\bo{A}\bo{x} = \bo{y}$ имеет точные решения, все они могут быть описаны как:
		
		\begin{equation}
			\bo{x} = \bo{A}^+\bo{y} + \bo{N}\bo{z}
		\end{equation}
		
		где $\bo{N}$ — базис в нулевом пространстве $\bo{A}$. Действительно, если $\bo{x}_p = \bo{A}^+\bo{y}$ является решением, то $\bo{x} = \bo{x}_p + \bo{N}\bo{z}$ тоже является решением, потому что $\bo{A}\bo{N}\bo{z} = 0$.
		
		\bigskip
		
		Однако нет гарантии, что $\bo{x}_p = \bo{A}^+\bo{y}$ является решением с \emph{нулевой невязкой}. Чтобы это проверить, убедимся, что невязка равна нулю:
		%
		\begin{align}
			\bo{y} - \bo{A}\bo{x}_p = 0, \ \ \text{или}
			\\
			(\bo{I} - \bo{A}\bo{A}^+)\bo{y} = 0
		\end{align}
		
		\bigskip
		
		Частное решение $\bo{x}_p = \bo{A}^+\bo{y}$ всегда \emph{ортогонально} решению из нулевого пространства $\bo{x}_n = \bo{N}\bo{z}$
		
	\end{flushleft}
\end{frame}




\begin{frame}{Литература}
	\begin{flushleft}
		
		\begin{itemize}
			
			\item \bref{https://ocw.mit.edu/courses/18-06sc-linear-algebra-fall-2011/0550c89b69c99e97dcbf52074e293308_MIT18_06SCF11_Ses3.8sum.pdf}{MIT, Left and right inverses; pseudoinverse.}
			
			
			\item \bref{https://faculty.math.illinois.edu/~mlavrov/docs/484-spring-2019/ch4lec4.pdf}{Minimum Norm Solutions, Math 484: Nonlinear Programming, Mikhail Lavrov}
			%https://faculty.math.illinois.edu/~mlavrov/courses/484-spring-2019.html
			
		\end{itemize}
		
		
	\end{flushleft}
\end{frame}




\begin{frame}{Упражнение}
	\begin{flushleft}
		
		\begin{itemize}
			\item Матрица $\bo{M}$ ортонормирована и квадратна, докажите, что $\bo{M}\T = \bo{M}^{-1}$.
			
			\item Найдите минимум $||\bo{A}\bo{x} - \bo{y}||_2$, когда столбцы $\bo{A}$ не являются линейно независимыми.
			
		\end{itemize}
		
		
		
	\end{flushleft}
\end{frame}




\myqrframe



\begin{frame}{Скалярное произведение и норма вектора}
	\begin{flushleft}
		
		Для двух векторов $\bo{x} = \begin{bmatrix}
			x_1 \\ x_2 \\ ... \\ x_n
		\end{bmatrix}$ и $\bo{y} = \begin{bmatrix}
			y_1 \\ y_2 \\ ... \\ y_n
		\end{bmatrix}$ их \emph{скалярное произведение} равно:
		
		\begin{equation}
			\bo{x} \cdot \bo{y} = x_1 y_1 + x_2 y_2 + ... + x_n y_n = \bo{x}\T \bo{y}
		\end{equation}
		
		\emph{2-норма} (также называемая евклидовой нормой) вектора определяется как:
		
		\begin{equation}
			||\bo{x}||_2 = \sqrt{\bo{x}\T \bo{x}} =  \sqrt{x_1 x_1 + x_2 x_2 + ... + x_n x_n}
		\end{equation}
		
		
	\end{flushleft}
\end{frame}





\begin{frame}{Производные}
	\begin{flushleft}
		
		Рассмотрим билинейную функцию $f(\bo{x}, \bo{y}) = \bo{y}\T \bo{M} \bo{x}$, где $\bo{x} \in \R^n$,  $\bo{y} \in \R^m$. Её производные (градиенты) равны:
		
		\begin{align}
			\frac{\partial f}{\partial \bo{x}} &= \bo{y}\T \bo{M} \in \R^{1 \times n}, 
			\\
			\frac{\partial f}{\partial \bo{y}} &= ( \bo{M} \bo{x})\T =  \bo{x}\T \bo{M}\T \in \R^{1 \times m}.
		\end{align}
		
		Для квадратичной формы $g(\bo{x}) = \bo{x}\T \bo{M} \bo{x}$ производная равна:
		
		\begin{align}
			\frac{\partial g}{\partial \bo{x}} &= \bo{x}\T \bo{M} + \bo{x}\T \bo{M}\T = \bo{x}\T (\bo{M} + \bo{M}\T).
		\end{align}
		
		
	\end{flushleft}
\end{frame}




\begin{frame}{Производные - пример}
	\begin{flushleft}
		
		Найдём частные производные функции:
		
		\begin{align}
			f(\bo{x}, \bo{y}) = \bo{y}\T \bo{M} \bo{x} = 
			\begin{bmatrix}
				y_1 \\ y_2
			\end{bmatrix}\T
			\begin{bmatrix}
				m_{11} & m_{12} \\
				m_{21} & m_{22}
			\end{bmatrix}
			\begin{bmatrix}
				x_1 \\ x_2
			\end{bmatrix} =
			\\
			= m_{11} x_1 y_1 + m_{12} x_2 y_1 + m_{21} x_1 y_2 + m_{22} x_2 y_2 .
		\end{align}
		%
		\begin{align}
			\frac{\partial f}{\partial x_1} = m_{11} y_1  + m_{21}  y_2,
			\ \ \
			\frac{\partial f}{\partial x_2} =  m_{12} y_1 + m_{22} y_2,
			%
			\\
			%
			\frac{\partial f}{\partial y_1} = m_{11} x_1 + m_{12} x_2,
			\ \ \
			\frac{\partial f}{\partial y_2} =  m_{21} x_1 + m_{22} x_2.
		\end{align}
		
		Используя это, мы можем составить градиенты:
		%
		\begin{align}
			\frac{\partial f}{\partial \bo{x}} &= 
			\begin{bmatrix}
				m_{11} y_1  + m_{21}  y_2 & m_{12} y_1 + m_{22} y_2
			\end{bmatrix}
			=
			\bo{y}\T \bo{M},
			\\
			\frac{\partial f}{\partial \bo{y}} &= 
			\begin{bmatrix}
				m_{11} x_1 + m_{12} x_2 & m_{21} x_1 + m_{22} x_2
			\end{bmatrix}
			=
			\bo{x}\T \bo{M}\T.
		\end{align}
		
		
	\end{flushleft}
\end{frame}



\begin{frame}{Приложение A, 1}
	\begin{flushleft}
		
		Псевдообратная матрица Мура-Пенроуза $\bo{M}^+$ должна удовлетворять следующим условиям:
		
		\begin{align}
			\bo{M}\bo{M}^+\bo{M} &= \bo{M} \\
			\bo{M}^+\bo{M} \bo{M}^+ &= \bo{M}^+ \\
			(\bo{M}\bo{M}^+)\T &= \bo{M}\bo{M}^+ \\
			(\bo{M}^+\bo{M})\T &= \bo{M}^+\bo{M}
		\end{align}
		
	\end{flushleft}
\end{frame}



\begin{frame}{Приложение A, 2}
	\begin{flushleft}
		
		Левая обратная матрица $\bo{A}^+ =\textcolor{myblue}{ (\bo{A}\T \bo{A})^{-1} \bo{A}\T}$ удовлетворяет условиям Мура-Пенроуза:
		
		\begin{align}
			\bo{A}\textcolor{myblue}{(\bo{A}\T \bo{A})^{-1} \bo{A}\T}\bo{A} &= \bo{A} 
			\\
			\textcolor{myblue}{(\bo{A}\T \bo{A})^{-1} \bo{A}\T} \bo{A} \textcolor{myblue}{(\bo{A}\T \bo{A})^{-1} \bo{A}\T} &= \textcolor{myblue}{(\bo{A}\T \bo{A})^{-1} \bo{A}\T} \\
			(\bo{A} \textcolor{myblue}{(\bo{A}\T \bo{A})^{-1} \bo{A}\T})\T &= \bo{A} \textcolor{myblue}{(\bo{A}\T \bo{A})^{-1} \bo{A}\T} \\
			( \textcolor{myblue}{(\bo{A}\T \bo{A})^{-1} \bo{A}\T}\bo{A})\T &= \textcolor{myblue}{(\bo{A}\T \bo{A})^{-1} \bo{A}\T} \bo{A}
		\end{align}
		
	\end{flushleft}
\end{frame}




\begin{frame}{Приложение B}
	\begin{flushleft}
		
		Мы можем сделать дополнительные наблюдения относительно формулы невязки $\bo{e} = \bo{y} - \bo{A}\bo{A}^+\bo{y} = (\bo{I} - \bo{A}\bo{A}^+)\bo{y} $:
		
		\begin{itemize}
			
			\item $\bo{A}\bo{A}^+\bo{y}$ является \emph{ортогональной проекцией} $\bo{y}$ на \emph{пространство столбцов} матрицы $\bo{A}$.
		
		\item $\bo{I} - \bo{A}\bo{A}^+$ является ортогональной проекцией $\bo{y}$ на \emph{левое нулевое пространство} матрицы $\bo{A}$.
	\end{itemize}
	
	Пространство столбцов — это "пространство всех возможных результатов" умножения на матрицу. Левое нулевое пространство — это пространство "всех результатов, которые матрица не может произвести". Оба пространства ортогональны: если $\bo{y}_1 \in \text{col}(\bo{A})$ и $\bo{y}_2 \in \text{left}(\bo{A})$, то $\bo{y}_1\T \bo{y}_2  = 0$.
	
\end{flushleft}
\end{frame}







\end{document}

